% Siteringstal for «dei ti namna» som ein rein tabell (ikkje eit diagram).
% Poenget er ikkje rangeringa, men kva talet faktisk tel.
\begin{table}[t]
\centering
\small
\begin{tabular}{@{}p{42mm}rp{58mm}@{}}
\toprule
\textsc{namn} & \textsc{siteringar} & \textsc{kva tyngda kviler på} \\
\midrule
Herbert Simon        & $\gg$\,500\,k & økonomi og kognisjon, ikkje design \\[2pt]
Donald Norman        & {\raise.17ex\hbox{$\scriptstyle\sim$}}\,170\,k & ei innflytelsesrik meining (\emph{affordans}) \\[2pt]
Klaus Krippendorff   & {\raise.17ex\hbox{$\scriptstyle\sim$}}\,120\,k & kommunikasjonsteori, lånt inn \\[2pt]
Nigel Cross          & {\raise.17ex\hbox{$\scriptstyle\sim$}}\,40\,k & «designerly ways of knowing» (program, ikkje resultat) \\[2pt]
Kees Dorst           & {\raise.17ex\hbox{$\scriptstyle\sim$}}\,21\,k & rammesetjing (framing) \\[2pt]
Victor Margolin      & {\raise.17ex\hbox{$\scriptstyle\sim$}}\,8\,k & designhistorie \\[2pt]
Richard Buchanan     & {\raise.17ex\hbox{$\scriptstyle\sim$}}\,7\,k & «wicked problems», lånt frå Rittel \\[2pt]
Bryan Lawson         & {\raise.17ex\hbox{$\scriptstyle\sim$}}\,7\,k & korleis designarar tenkjer \\[2pt]
Ezio Manzini         & {\raise.17ex\hbox{$\scriptstyle\sim$}}\,6\,k & sosial innovasjon \\[2pt]
Nelson \& Stolterman & {\raise.17ex\hbox{$\scriptstyle\sim$}}\,5\,k & design som eigen kunnskapsform \\
\bottomrule
\end{tabular}
\caption{\textsc{Siteringar for «dei ti namna».} Omtrentlege tal (Google Scholar, storleiksorden). Poenget er ikkje rangeringa, men kva talet tel: dei mest siterte er det oftast for arbeid utanfor design, eller for éi særskilt innflytelsesrik meining --- ikkje for kumulativ, etterprøvd kunnskap som har bygd seg opp i faget.}
\end{table}
